% The cos(2 phi) transmon (KITE) of Roverch et al., drawn as in the source: a
% hexagonal loop with a junction on each vertical side and an inductor on each diagonal,
% threaded by Phi_ext, shunted by C_S between the top node and ground. Only the
% shunted node carries capacitance to ground, so the two modes in which a junction's
% terminals move together are frozen. Reference frame: varphi (compact, across C_S),
% phi_Sigma and phi_Delta (sum and difference of the two upper inductor phases).
% Compile standalone with the draw-circuit driver; \input into main.tex via the
% standalone package. The josephson bipole comes from jj_bipoles.tex, which main.tex also loads.
\documentclass[border=3pt]{standalone}
\usepackage{tikz}
\usepackage{circuitikz}
% Shared circuitikz bipole declarations, copied from the draw-circuit skill template.
% Loaded in the preamble of main.tex and of each standalone figure in figs/.
% Boxed Josephson junction: the co-located E_J || C_J branch with zero loop area.
\ctikzset{
    bipoles/josephson/height/.initial = 0.35,
    bipoles/josephson/width/.initial  = 0.35,
}
\makeatletter
\pgfcircdeclarebipolescaled{misc}
{}
{\ctikzvalof{bipoles/josephson/height}}
{josephson}
{\ctikzvalof{bipoles/josephson/height}}
{\ctikzvalof{bipoles/josephson/width}}
{
    \pgf@circ@setlinewidth{bipoles}{\pgfstartlinewidth}
    \pgfscope
        \pgfpathrectanglecorners{\pgfpoint{\pgf@circ@res@left}{\pgf@circ@res@down}}{\pgfpoint{\pgf@circ@res@right}{\pgf@circ@res@up}}
        \pgfsetfillcolor{white}
        \pgfusepath{fill,stroke}
        \pgfpathmoveto{\pgfpoint{\pgf@circ@res@left}{\pgf@circ@res@down}}
        \pgfpathlineto{\pgfpoint{\pgf@circ@res@right}{\pgf@circ@res@up}}
        \pgfpathmoveto{\pgfpoint{\pgf@circ@res@left}{\pgf@circ@res@up}}
        \pgfpathlineto{\pgfpoint{\pgf@circ@res@right}{\pgf@circ@res@down}}
        \pgfusepath{stroke}
    \endpgfscope
}
\def\pgf@circ@josephson@path#1{\pgf@circ@bipole@path{josephson}{#1}}
\tikzset{josephson/.style = {\circuitikzbasekey, /tikz/to path=\pgf@circ@josephson@path, l=#1}}

\makeatother

\begin{document}
\begin{circuitikz}[line width=0.7pt, font=\small]
    \ctikzset{label/align = straight, bipoles/thickness=1, capacitors/scale=0.55, inductors/scale=0.9, inductors/coils=4, misc/scale=0.8}
    \tikzset{lbl/.style={font=\small, inner sep=0.5pt}}
    \coordinate (T) at (0,1.7);
    \coordinate (B) at (0,-1.7);
    \coordinate (UL) at (-1.472,0.85);
    \coordinate (LL) at (-1.472,-0.85);
    \coordinate (UR) at (1.472,0.85);
    \coordinate (LR) at (1.472,-0.85);
    \coordinate (CT) at (-2.2,1.7);
    \coordinate (CB) at (-2.2,-1.7);
    % left arm
    \draw (T) to[L, l=$L$] (UL);
    \draw (UL) to[josephson] (LL);
    \draw (LL) to[L, l=$L$] (B);
    % right arm
    \draw (T) to[L, l_=$L$] (UR);
    \draw (UR) to[josephson] (LR);
    \draw (LR) to[L, l_=$L$] (B);
    % shunt on the left and ground
    \draw (T) -- (CT);
    \draw (CT) to[C, l_=$C_S$] (CB);
    \draw (CB) -- (B);
    \draw (B) node[ground]{};
    % junction labels inside the loop, flux at the center
    \node[lbl] at (-0.98,0) {$E_{J1}$};
    \node[lbl] at (0.98,0) {$E_{J2}$};
    \node at (0,0) {$\odot\,\Phi_{\rm ext}$};
\end{circuitikz}
\end{document}
